% =====================================================================
% SECTION 08 — Discussion: Limitations
%
% Slot: Append to or replace the existing Limitations subsection in
% the paper's Discussion / Conclusion section.
% Suggested label: \label{sec:limitations_gastown}
%
% This section covers ONLY the Gastown-integration-specific limitations.
% The existing paper's limitations section (synthetic corpus scope,
% iteration-budget dependence, etc.) should be retained as-is.
%
% Word target: 400--800 words.
% =====================================================================

\subsection{Limitations of the Gastown Integration and Schema-Gap Claims}
\label{sec:limitations_gastown}

\paragraph{Corpus scope: bead-tracker only.}
The Phase 1A baseline ($3{,}667$ events, $31$ agents) derives exclusively
from Gastown's committed bead-tracker archive
(\texttt{.beads/backup/events.jsonl}). The runtime session feed
(\texttt{\textasciitilde/gt/.events.jsonl}), which carries patrol events,
escalations, slings, handoffs, and merge outcomes, is generated at runtime
and not committed to the repository. It is therefore absent from our
corpus. The runtime feed would provide substantially richer behavioral
signal: the $8$ silent brains were specifically designed for event types
that the runtime feed carries (\texttt{patrol\_started},
\texttt{escalation\_sent}, \texttt{sling}, \texttt{hook},
\texttt{handoff}) but the bead-tracker archive does not. The all-CLEAR
finding is a property of the available corpus, not a claim about Gastown
agent safety at runtime.

\paragraph{Refinery gate: 3-brain PoC scope.}
The Refinery gate demonstration (\S\ref{sec:gastown:protocol}, Phase 2C)
uses ContrastiveBrain, IdentityBrain, and TemporalBrain---the three brains
that produced non-zero scores on the bead-tracker corpus. Note that
ContrastiveBrain's scores in the gate reflect a saturation artifact
(\S\ref{sec:gastown:diagnosis}); the gate's calibration thresholds are
correct for the Phase 1A corpus but should be revisited when a substrate-adaptive
reference distribution is available. The $8$ silent brains
are excluded from the gate calibration. Gate thresholds (warn at the
$90$th percentile, pass at the $95$th percentile of the CLEAR-band
distribution) are calibrated on $31$ agents, all CLEAR, with no
known-adversarial agents in the calibration set. A corpus that includes
known-suspicious agent behavior would produce better-calibrated thresholds;
the current calibration is explicitly a starting point, not a production
recommendation.

\paragraph{Synthetic fixtures for silent-brain validation.}
Phase 2B validated that the Mayor tap infrastructure and Witness third gate
can produce behavioral signal from previously-silent brains---but this
validation was performed against synthetic fixtures, not live Gastown
events. The synthetic patterns (zombie detection, escalation streams,
injection attempts) are designed to match Gastown's documented
\texttt{GtEvent} schema, but they are not drawn from a real Gastown
deployment. Results from synthetic fixture replay should be treated as
wiring validation, not empirical detection claims.

\paragraph{Gas City migration risk.}
The evidence matrix (\S\ref{sec:abos:problem}) is based on Gastown at
commit HEAD (April 20, 2026). Kilo Code has announced
Gas City~\cite{kilocode2026gasCity}---a managed cloud platform for
Gastown-style multi-agent workflows. If Gas City ships a native behavioral
observability layer, some ABOS fields may be rendered redundant or may
conflict with a proprietary schema. The ABOS specification was designed to
be backward-compatible; we recommend that any Gas City behavioral telemetry
adopt ABOS fields rather than parallel naming to avoid fragmentation.

\paragraph{Ensemble overhead not benchmarked on production loads.}
The $13$-brain ensemble has not been benchmarked on production-scale
orchestration loads. Phase 2C's Refinery gate uses a $3$-brain subset
with a $30$-second timeout per merge event---a conservative but unbenchmarked
configuration. At high throughput (e.g., $> 100$ merge events per hour),
the gate may become a bottleneck. The ABOS MVS overhead estimate of
$< 2$\,ms per event (\S\ref{sec:abos:mvs}) is derived from operation type
analysis, not from measured wall-clock benchmarks on a real orchestrator.
Production deployments should benchmark the ensemble under their specific
event rates before enabling blocking-mode gates.

\paragraph{Cross-substrate confidence bands.}
The cross-substrate activation counts (Table~\ref{tab:cross_substrate}) are
reported with explicit confidence bands. ContrastiveBrain is excluded from
the comparison as a saturation-prone instrument; see \S\ref{sec:gastown:diagnosis}
for the full diagnosis. The OpenHands result ($1/12$ instrumentally-valid
brains, LOW confidence) is based on $52$ events from two SWE-bench tasks
with absent timestamps in the test fixture; it is not representative of
a production OpenHands session. The Claude Code result ($5/12$ brains,
MEDIUM confidence) is based on $500$ records from a single session; it
has not been validated across a diverse sample of Claude Code sessions.
Both results should be treated as directional indicators of the
schema-gradient hypothesis, not as definitive cross-substrate comparisons.

\paragraph{Requested API additions block Phase C deployment.}
The Behavioral Warrant's Phase C deployment (live gating) requires two
additions to Gastown's \texttt{stamps} table that are not in the current
codebase: (1) support for a non-rig \texttt{author} value
(\texttt{"sybilcore-warrant"} as a system-actor identity), and (2) an
Ed25519 \texttt{signature} field and associated key registry in Wasteland
governance records. Until these additions are prioritized by Gastown
governance, the Warrant operates in observe-only mode against a shadow
table, and the formal portability and signing semantics of
\S\ref{sec:warrant:portability} remain unimplemented.